\section{Path geometry and exact projection}\label{app:path-geometry}
This appendix supplies the statements used in Section~\ref{subsec:follower-path}.
The path shadow is an established construction. The telescoping identity gives
a second exposure formula tailored to the quadratic follower, and the final
projection result applies to the narrower class of equal-gain strips.

\subsection{The sparse Klee--Minty shadow and its messages}
\label{app:klee-shadow}
Every coordinate of \(P_n(\gamma)\) lies in \([0,1]\). Since \(\gamma<1/2\),
the two bounds on a coordinate cannot be simultaneously tight. A vertex has
\(n\) independent tight rows, hence exactly one bound per coordinate. Conversely,
every endpoint selection is nonsingular triangular. The vertices are therefore
\begin{equation}\label{eq:klee-vertices}
 v_j(d)=d_j+(1-2d_j)\gamma v_{j-1}(d),\quad v_0=0,
 \qquad d\in\{0,1\}^n.
\end{equation}
Two vertices are adjacent when their bit vectors differ in one coordinate.
Their terminal values are distinct: the two final choices lie in the disjoint
intervals \([0,\gamma]\) and \([1-\gamma,1]\); within either interval the
terminal map is injective in the previous coordinate, so induction applies.

\begin{lemma}[Sparse shadow, after G\"artner et al.]\label{lem:klee-shadow}
Let \(\gamma=1/4\), \(\widehat c_j=\gamma^{3(n-j)}\) for \(j<n\), and
\(\widehat c_n=0\). Every vertex \(v(d)\) uniquely maximizes
\(\widehat c\trans z+\lambda z_n\) over \(P_n(\gamma)\) at the price
\[
 \lambda(d)=-\sum_{j=0}^{n-1}
       \left(\prod_{k=j+1}^n(1-2d_k)\right)\gamma^{2(n-j)}
       \in(-1/15,1/15).
\]
\end{lemma}
\begin{proof}
We reproduce the edge calculation in \citet[Definition~11 and Lemma~12]{GartnerEtAl2013}
to fix the exact version used here. Put
\(p_i^j=\prod_{k=i}^j(1-2d_k)\), with an empty product equal to one.
If bit \(\ell\) is flipped, the difference in coordinate \(j\ge\ell\) is
\(p_{\ell+1}^j\gamma^{j-\ell}q_\ell\), where
\(q_\ell=(1-2d_\ell)(1-2\gamma v_{\ell-1})\); earlier coordinates do not
change. Divide the objective difference along this edge by \(q_\ell\).
The contributions with \(j\ge\ell\) cancel between \(\widehat c\) and
\(\lambda(d)e_n\), leaving
\[
 -\sum_{j=0}^{\ell-1}p_{j+1}^n p_{\ell+1}^n
                  \gamma^{3n-\ell-2j}.
\]
The lowest-power coefficient, at \(j=\ell-1\), is \(-(1-2d_\ell)\).
All other coefficients have magnitude one and their powers increase by two;
their total magnitude relative to the leading term is less than
\(\gamma^2/(1-\gamma^2)<1\). Multiplying by \(q_\ell\) gives a strictly
negative edge difference. The feasible direction cone at a simple polytope
vertex is generated by its incident edge directions (solve in its independent
active-row coordinates); hence negativity on every edge implies unique global
maximality. Finally
\(|\lambda(d)|\le\sum_{h=1}^n\gamma^{2h}<\gamma^2/(1-\gamma^2)=1/15\).
\end{proof}

\begin{proof}[Proof of Proposition~\ref{prop:path-shadow-boundary}]
Let \(V_n(\lambda)=\max_{z\in P_n}\widehat c\trans z+\lambda z_n\).
Each of its \(2^n\) vertex lines is uniquely maximal at the price in
Lemma~\ref{lem:klee-shadow}, and remains so on an open neighborhood. A maximum
of affine functions cannot leave and later return to the same strictly active
line, since that line's dominance set is an interval. Thus all \(2^n\) lines
give distinct maximal affine pieces.
Pointwise evaluation does not require their construction. Eliminate the last
coordinate, and continue backwards using
\[
 a_n=\lambda,\qquad a_j=\widehat c_j-\gamma|a_{j+1}|\ (j=n-1,\ldots,1),
 \qquad V_n(\lambda)=\sum_{j=1}^n\max(0,a_j).
\]
A nonnegative coefficient selects \(z_j=1-\gamma z_{j-1}\), contributing its
constant and subtracting \(\gamma a_j\) from the previous coefficient. A
negative coefficient selects \(z_j=\gamma z_{j-1}\), contributing no constant
and adding \(\gamma a_j\). At zero either choice is valid. This proves the
recurrence. It uses \(O(n)\) rational operations, with bit lengths growing
polynomially in \(n\) and the input price length.
The terminal-state message
\(M_n(t)=\max\{\widehat c\trans z:z\in P_n,z_n=t\}\) is the upper boundary
of the planar projection \((z_n,\widehat c\trans z)\). Every projected vertex
is uniquely exposed by a direction with coefficient one on the second
coordinate, so all belong to this upper boundary. Their terminal coordinates
are distinct, and no three consecutive projected vertices are collinear
(otherwise the middle one could not be uniquely exposed). The extreme terminal
values zero and one occur. Hence \(M_n\) has \(2^n-1\) maximal affine intervals.

Consider the bounded epigraph
\[
 E_n=\{(\lambda,v):-1/15\le\lambda\le1/15,\quad
                         V_n(\lambda)\le v\le1\}.
\]
Here \(V_n<1\), since \(\sum_j\widehat c_j<1/63\) and \(|\lambda|\le1/15\).
Suppose a quantifier-free Boolean formula in these two variables uses the
nonzero polynomials \(P_1,\ldots,P_s\) to describe \(E_n\).
At every point on an open linear piece of its lower boundary, some polynomial
must vanish: if all were nonzero, their signs and thus the formula's truth
would be constant in a neighborhood. The product \(\prod_iP_i\) therefore
vanishes on each of \(2^n\) distinct open line segments. Its restriction to
each supporting line is a univariate polynomial vanishing on an interval, so
that line's defining linear polynomial is a factor. Distinct line factors are
coprime; their product divides \(\prod_iP_i\). Thus
\(\sum_i\deg P_i\ge2^n\). Identically zero predicates can be removed from
the formula in advance. This lower bound explicitly excludes auxiliary
variables and quantified or circuit representations.
\end{proof}

\subsection{A vertex polynomial, a slab, and padding}
\begin{lemma}[Telescoping and parabolic exposure]\label{lem:klee-telescope}
For \(P_n(\gamma)\), let
\[
 L(z)=z_n-(1-\gamma)\sum_{i<n}\gamma^{2(n-i)-1}z_i,
 \qquad F(z)=L(z)-z_n^2.
\]
Then
\begin{equation}\label{eq:klee-telescoping}
 F(z)=\sum_{j=1}^n\gamma^{2(n-j)}
   (z_j-\gamma z_{j-1})(1-\gamma z_{j-1}-z_j)\ge0.
\end{equation}
Its zeros in \(P_n(\gamma)\) are exactly the vertices. Every vertex with
terminal value \(t\) uniquely maximizes \(2tz_n-L(z)\). The Hessian of \(F\)
is \(-2e_ne_n\trans\).
\end{lemma}
\begin{proof}
The product in the \(j\)-th term expands to
\(z_j-z_j^2-\gamma z_{j-1}+\gamma^2z_{j-1}^2\).
The weighted quadratic terms telescope to \(-z_n^2\); the interior linear
coefficient is
\(\gamma^{2(n-j)}-\gamma^{2(n-j)-1}=-(1-\gamma)\gamma^{2(n-j)-1}\).
This proves the identity. Both product factors are nonnegative on the
polytope, and neither pair can simultaneously vanish. Equality therefore
selects one endpoint per coordinate, precisely \eqref{eq:klee-vertices}.
For any \(z\in P_n\),
\[
 2tz_n-L(z)\le2tz_n-z_n^2=t^2-(z_n-t)^2\le t^2.
\]
Equality forces \(F=0\) and \(z_n=t\), hence the unique vertex with that
terminal value. The Hessian follows directly from the definition.
\end{proof}

\begin{lemma}[Subset Sum slab and fixed local coefficients]\label{lem:klee-padding}
For positive integers \(a_i\), \(W=\sum_i a_i\), and \(0\le T\le W\), the
system \(z\in P_n(1/(8W))\), \(F(z)\le0\), and
\(T-1/4\le\sum_i a_i z_i\le T+1/4\) is feasible precisely for a positive
Subset Sum instance. The same equivalence holds with \(\gamma=1/4\), additional
unary padding bounds, and polynomially increased dimension. Its feasible set
without \(F\le0\) is nonempty and compact.
\end{lemma}
\begin{proof}
For the unpadded system, a vertex's coordinate differs from its bit by at most
\(\gamma\). A satisfying bit pattern therefore gives slab error at most
\(\gamma W=1/8\). Conversely, \(F\le0\) forces a vertex, and its integer
subset sum differs from \(T\) by at most \(1/4+1/8=3/8<1\), so it equals
\(T\).
For fixed \(\gamma=1/4\), take the least integer \(r\ge0\) such that
\(4^{-(r+1)}\le1/(8W)\). Set \(N=n+(n-1)r\), with free positions
\(j_i=1+(i-1)(r+1)\), and require \(z_j\le1/2\) at every other position.
At an original cube vertex, an upper endpoint is at least \(3/4\). Thus all
padding coordinates take their lower endpoints. Between successive free
positions there are \(r\) multiplications by \(1/4\), and
\[
 |z_{j_i}-d_{j_i}|\le4^{-(r+1)}\quad(i>1),\qquad z_1=d_1.
\]
Every free bit pattern extends by zero padding bits to an original cube
vertex satisfying all unary bounds. Apply the slab only to
\(\sum_i a_i z_{j_i}\). The same \(1/8\) and \(3/8\) estimates prove the
equivalence. Here \(r=O(\log W)\), so dimension and bit lengths are polynomial.
The local rows use coefficients \(0,\pm1,\pm1/4\) and right-hand sides
\(0,1/2,1\); only the aggregate and vertex polynomial retain variable binary
data. For nonemptiness without the vertex row, the zero vertex has sum zero,
and the vertex with every free bit one has sum at least \(W-1/8\).
Their segment reaches every integer \(T<W\); for \(T=W\) the latter vertex
already meets the slab. Compactness follows from the unit box.
\end{proof}

\subsection{Equal-gain strips: a polynomial projection}
\label{app:strip-projection}
\begin{proposition}[Paths, trees, and fixed total cycle rank]
\label{prop:strip-projection}
Fix the dimension of a parameter \(t\) and the number of retained state
variables. Suppose scalar states have finite bounds, with rational polynomial
coefficients in \(t\), and each directed edge \(u\to v\) imposes
\begin{equation}\label{eq:affine-strip}
 a_e(t)x_u+b_e(t)\le x_v\le a_e(t)x_u+c_e(t).
\end{equation}
If the underlying graph is a forest, one can eliminate all nonretained states
with polynomial description size, numerical degree and bit complexity.
The same holds for fixed total cycle rank. Signs of gains, including zero,
are unrestricted. Additional objectives and constraints may depend on the
parameters and retained states only. For fixed remaining dimension, exact
feasibility, finite-infimum computation and attainment decisions consequently
have polynomial bit complexity under dense polynomial encoding. Witnesses
can be recovered in the same real number field as the retained solution.
\end{proposition}
\begin{proof}
First consider a path with nonzero gains \(a_i\), indexed so that its edges are
\(i-1\to i\). On a realizable sign condition of the input gains put
\(A_0=1,A_i=\prod_{h=1}^i a_h\), and \(y_i=x_i/A_i\).
The node bounds become \(L_i\le y_i\le U_i\), with endpoints interchanged
when \(A_i<0\). The transition becomes
\(\alpha_i\le y_i-y_{i-1}\le\beta_i\), obtained from \(b_i/A_i,c_i/A_i\)
with the same sign rule. Require each interval to be consistent. Interpret
these as directed difference constraints, with edge length \(\beta_i\) from
\(i-1\) to \(i\) and \(-\alpha_i\) in the reverse direction. Every two-edge
cycle has nonnegative length. Removing backtracks shows that the unique
simple path between two vertices is a shortest path. Denote its length by
\(d(j,i)\), including \(d(i,i)=0\). Thus feasibility is exactly
\begin{equation}\label{eq:strip-all-pairs}
 L_i\le U_j+d(j,i)\qquad\text{for every pair }i,j.
\end{equation}
Necessity follows by summing along a path. For sufficiency, set
\begin{equation}\label{eq:strip-recovery}
 y_i=\min_j\{U_j+d(j,i)\}.
\end{equation}
The all-pairs inequalities give the lower bounds, the term \(j=i\) gives the
upper bounds, and the shortest-path triangle inequality gives every directed
edge inequality. To retain a specified state, add its prescribed value as
both a lower and an upper bound, while keeping its original bounds. This
makes \eqref{eq:strip-all-pairs} an exact projection onto the parameters and
retained values, rather than only a feasibility test.

If a gain is zero, its original row imposes only \(b_e\le x_v\le c_e\).
Delete the edge and add these bounds to its original head \(v\); its original
direction matters here. Restart the products in each remaining component.
Keep multiple lower and upper candidates at a node without selecting a
largest or smallest one. Require every lower candidate to be at most every
propagated upper candidate; in \eqref{eq:strip-recovery}, minimize over all
upper candidates. This covers isolated nodes, consecutive zero gains,
inconsistent bounds and lower-dimensional parameter sign conditions.

For an arbitrarily oriented tree with nonzero gains, choose a root and factors
\(A_v\) satisfying \(A_v=a_e A_u\) on each original edge \(u\to v\).
Traversing against an edge uses its reciprocal gain. The normalization
\(y_v=x_v/A_v\) again makes each strip a bounded difference interval.
The unique-simple-path proof, all-pairs conditions, retained-value bounds and
minimum recovery all remain valid. Zero edges are removed before this
normalization as above. With \(O(N)\) bound candidates there are \(O(N^2)\)
all-pairs conditions. Removing \(c\) nonforest edges from a graph retains at
most \(2c\) further endpoint states. Project the forest with these endpoints
retained and then reinsert the original removed rows. For fixed \emph{total}
cycle rank \(c\), the remaining dimension stays fixed. A bound on the cycle
rank of each of many blocks would not suffice.

It remains to check the symbolic and bit bounds. Realizable sign conditions
of the input gains can be enumerated in polynomial time in fixed parameter
dimension, retaining zeros. Within each, the normalization is a rational
function with a denominator dividing a product of input nonzero gains; tree
reciprocals cause no iterative algebraic extensions. Prefix products, path
sums and all-pairs conditions use only polynomially many such factors.
One can use a common product of all nonzero gains, with additional powers if
needed, to clear denominators of known sign, or multiply by its square and
retain its nonvanishing condition. The degrees are \(O(N\delta)\) up to a
fixed polynomial factor when input degree is \(\delta\), and the number of
monomials is polynomial in this degree because the parameter dimension is
fixed. Exact multiplication and addition therefore yield polynomial coefficient
bit length and total description size. The fixed number of retained state
variables only adds a fixed number of polynomial coordinates. The resulting
union of projected sign conditions is a fixed-dimensional semialgebraic
formula, to which the operations of Section~\ref{subsec:tools} apply.
Finally, \eqref{eq:strip-recovery} selects one of finitely many rational functions
of the retained solution; comparisons and rational arithmetic stay in its
common field. Multiplication by \(A_i\) recovers each original state.

Finite state intervals alone do not make the parameter domain compact. The
conclusion therefore distinguishes a finite infimum from an attained minimum;
an optimizer is returned when attained, and unboundedness can also be detected.
Under a closed compact feasible-domain assumption a continuous objective
attains its minimum. An objective or additional row involving all eliminated
states is outside the projection claim unless independently expressible in the
retained variables. The common gain on both sides of
\eqref{eq:affine-strip} is essential.
\end{proof}
